% MOVED OUT OF THE MAIN TEXT on 28 Oct for the ten-page limit (the author's decision;
% notes/22_page_cut_plan.md). This file is now input by appendix.tex, not by main.tex,
% which is why its heading is \section and its main-text references are numbers.
% Privacy-first section 4: recovery near each run's best saved held-out loss.
% Budget 1.75 pages including Table~\ref{tab:cost} and Figure~\ref{fig:cost}.
% Every number here is from the single-basis re-measurement of 19-23 Oct; see notes 3.58-3.59.
\section{Sequence recovery near the best saved held-out loss}
\label{sec:leak}

The question a practitioner faces is whether leakage appears at a checkpoint they would choose,
not whether a model memorizes after being trained past the point of usefulness. We
therefore pair every leakage rate with its held-out loss cost relative to the
best checkpoint that run itself produced. The decoding protocol and the number of record insertions jointly determine whether reproduction
occurs (Figure~\ref{fig:overview}b, upper half).
The second quantity makes the first operationally meaningful,
because it states how much predictive loss accompanies the measured recovery.

\paragraph{The cheapest leak costs seven thousandths of a bit.}
Table~\ref{tab:cost} orders the released checkpoints by that price. DNABERT-2 reproduces 30\% of
planted records at \textbf{+0.007\,\bits} above its own optimum. Nucleotide Transformer gives up
\textbf{+0.03\,\bits} to leak 8--14\% across three seeds. These checkpoints remain close to the best saved held-out loss, despite measurable recovery.

The never-trained control remains at zero throughout these evaluations. Records the model never saw are reproduced 0\% of the time in every
cell, every tier and every seed, and the same decoder scores the never-trained window at the foot of
each block on the same host, so the zero is a property of the measurement and not of a
separate baseline.

\begin{table}[tbp]
\centering
\caption{Released masked checkpoints on \emph{E.~coli}, records planted $16\times$, three seeds where shown. NT uses 50 records per tier;
the DNABERT-2 round-30 prefix and search cells use 10 and 20 records, respectively. \emph{Cost} is held-out \bits{} above that run's own best saved
checkpoint. \emph{Standard} is the prefix-conditioned protocol the field uses; \emph{ours} is the
re-masking decoder of \S\ref{sec:measure}. The DNABERT-2 rates are 3/10 and 0/20,
so that row compares protocols with different record coverage.}
\label{tab:cost}
\begin{tabular}{llccc}
\toprule
checkpoint & corpus & cost & standard & ours \\
\midrule
DNABERT-2 117M      & \emph{E.~coli} & \textbf{+0.007} & \textbf{30\%}      & 0\% \\
NT 250M (3 seeds)   & \emph{E.~coli} & \textbf{+0.03}  & 4 / 8 / 4\%        & \textbf{8 / 8 / 14\%} \\
NT 250M (3 seeds)   & \emph{E.~coli} & +0.44           & 26 / 26 / 30\%     & \textbf{66 / 68 / 76\%} \\
NT 500M             & \emph{E.~coli} & +0.46           & 36\%               & \textbf{78\%} \\
\midrule
never-trained control & --- & --- & 0\% & 0\% \\
\bottomrule
\end{tabular}
\end{table}

\paragraph{The two protocols disagree, and neither dominates.}
Comparisons across Table~\ref{tab:cost} show why this measurement cannot rely on a single extraction
test. On fixed-width non-overlapping tokenizers the standard protocol under-reports by
roughly a factor of two: at +0.44\,\bits{} it recovers 26--30\% where ours recovers 66--76\%. On
DNABERT-2 the order reverses completely. The standard protocol recovers 30\% and ours recovers
nothing at all. \S\ref{sec:measure} explains this protocol dependence. For each model, we interpret the reported exposure
using the stronger tested protocol. These measurements therefore provide lower bounds on what an attacker could obtain and not artifacts of one decoder's blind spot.

\paragraph{A single copy of a human-variant-derived record can be recovered.}
Uniform random canaries provide a controlled test but omit the sequence structure
of genomic inputs. We therefore extend the measurement to chromosome-22 reference
windows carrying observed rare alleles, testing whether recovery also occurs beyond
uniform random strings. These are the variant-derived records of \S\ref{sec:measure}. This is the original
mixed collection: 145 of 400 windows are chromosome-compatible at the audited sites,
20 are within-person genotype composites, and 235 are unobserved composites
(\S\ref{app:haplo}). Compatibility remains qualified at indel and multiallelic sites.
Recovery therefore demonstrates reproduction of a planted variant-derived record; it
does not establish full-base authenticity or individual linkage.

\begin{table}[tbp]
\centering
\caption{Nucleotide Transformer fine-tuned on human chromosome~22 with constructed rare-variant records
as canaries, including the mixed categories described above. \emph{1 copy} means one record appearing once in the
fine-tuning set. The second column is held-out \bits{} at that checkpoint. Round 50 is three seeds,
decoded in one batch; the earlier rounds are the first seed only, and are there to show the
trajectory rather than to carry an interval.}
\label{tab:human}
\begin{tabular}{lccccc}
\toprule
& & \multicolumn{2}{c}{16 copies} & \multicolumn{2}{c}{1 copy} \\
\cmidrule(lr){3-4}\cmidrule(lr){5-6}
checkpoint & held-out \bits & standard & ours & standard & ours \\
\midrule
NT 250M, round 15          & 1.411         & 6\%      & 6\%                & 0\%    & 0\% \\
NT 250M, round 30          & 1.501         & 10\%     & 28\%               & 0\%    & 0\% \\
NT 250M, round 50, 3 seeds & 1.632--1.648  & 16--18\% & \textbf{48--54\%}  & 0--2\% & \textbf{8\%} \\
NT 500M, round 15          & 1.372         & 16\%     & 14\%               & 0\%    & 0\% \\
NT 500M, round 50, 3 seeds & 1.668--1.788  & 24--50\% & \textbf{58--72\%}  & 0\%    & \textbf{12--18\%} \\
\midrule
never-trained control & --- & 0\% & 0\% & 0\% & 0\% \\
\bottomrule
\end{tabular}
\end{table}

Table~\ref{tab:human} presents the central result of this section. A variant-derived
record, present exactly once in the fine-tuning data, is recovered exactly by the 250M model in 8\% of cases in \emph{every one} of three seeds, and by the 500M in
12--18\%, between one case in twelve and one in six, after our fine-tuning of publicly downloadable starting weights, with the
matched never-trained control at zero in all of them. The repeated tier
behaves as the \emph{E.~coli} arm does, so the single-copy column is not an artifact of a canary
appearing often enough to be learned as vocabulary.

\paragraph{Conditional prediction and measured recovery.}
For DNABERT-2 on chromosome~22, single-copy records at the final checkpoint have
93--95\% nucleotide-level LOO accuracy and a power summary of 0.21--0.39 across
three seeds, while both tested decoders recover none of their hidden targets
(Table~\ref{tab:loo}). LOO supplies every other true target nucleotide; these are
conditional-prediction diagnostics, not bounds or attainable whole-record recovery
rates (\S\ref{sec:measure}).
Separately, for the \emph{E.~coli} sixteen-copy condition in Table~\ref{tab:cost},
DNABERT-2's prefix protocol recovers 30\% whereas the tested re-masking search
recovers 0\%. We report the stronger measured protocol for that condition.
Neither search failure alone establishes an absence of record-specific learning.

Table~\ref{tab:loo} compares the listed released-model checkpoints on chromosome~22, because the
same caution distinguishes conditional prediction with true target context from whole-record
decoding without that context. Read down the single-copy rows: the overlapping DNABERT family is the
case with the smallest power summary, $0.01$--$0.02$, alongside zero tested recovery.
Elsewhere the power summary is 0.14 to 0.41 while measured recovery is 0--18\%.
These quantities describe different tasks: the probe receives the other true target
positions, whereas the decoder must infer them. No power value here demonstrates that
a better decoder can attain it or excludes recovery by another protocol.

\begin{table}[tbp]
\centering
\caption{Conditional prediction and tested recovery on chromosome~22, 50 records per tier,
three seeds where shown. LOO supplies all other true target positions. \emph{Power diag.}
is the mean of per-record $p_i^{s_i}$, not a recovery bound or the power of mean accuracy.
NT probes tokens; DNABERT-2 and overlapping DNABERT probe 48 nucleotides. Control powers
are $0.000$--$0.003$. We rescored the overlapping-family prefix column at all fifty
records (\S\ref{app:prov}). Protocols and budgets are specified in \S\ref{sec:decoders}.}
\label{tab:loo}
\begin{tabular}{llcccc}
\toprule
checkpoint & copies & LOO top-1 & power diag. & standard & ours \\
\midrule
NT 250M, round 50 (3 seeds) & 16 & 0.96--0.97 & 0.81--0.85 & 16--18\% & \textbf{48--54\%} \\
                       & 1  & 0.57--0.66 & 0.14--0.23 & 0--2\%   & \textbf{8\%}  \\
NT 500M, round 50 (3 seeds) & 16 & 0.96--0.99 & 0.84--0.92 & 24--50\% & \textbf{58--72\%} \\
                       & 1  & 0.67--0.80 & 0.24--0.41 & 0\%      & \textbf{12--18\%} \\
DNABERT-2, final (3 seeds) & 16 & 0.99--1.00 & 0.85--0.91 & \textbf{40--44\%} & 0\% \\
                       & 1  & 0.93--0.95  & 0.21--0.39  & 0\%     & 0\% \\
DNABERT $k=3$--6, round 120 & 16 & 0.99--1.00 & 0.70--0.87 & 0--50\%  & 0--4\% \\
                       & 1  & 0.61--0.66  & 0.01--0.02  & 0\%     & 0\% \\
\midrule
never-trained control  & 0  & 0.07--0.51  & 0.000--0.003 & 0\%     & 0\% \\
\bottomrule
\end{tabular}
\end{table}
